\chapter{Egison早巡り}\label{quick}

本章は，パターンマッチ指向プログラミングのためのEgisonの機能を一通り紹介する．

\section{\lstinline{matchAll}式によるパターンマッチ}

パターンマッチを記述するためのいくつかの構文をEgisonは提供している．
そのうちもっとも基本的な構文は\lstinline{matchAll}\index{matchAll@\lstinline{matchAll}}式である．

\begin{lstlisting}[language=egison]
matchAll [1,2,3] as list something with
| $x :: $ts -> (x, ts)
-- [(1,[2,3])]
\end{lstlisting}

\lstinline{matchAll}式は，ターゲット（上記の場合， \lstinline{[1,2,3]} ），マッチャー（上記の場合，\lstinline{list something}），マッチ節（上記の場合，\lstinline{$x :: $xs -> (x,xs)}）から構成される．
マッチ節は，パターン（上記の場合，\lstinline{$x :: $xs}）とボディ（上記の場合，\lstinline{(x,xs)}）からなる．
\lstinline{matchAll}式は，既存のプログラミング言語のマッチ式と同様に，ターゲットとパターンのパターンマッチを試みて，もしマッチしたらそのマッチ節のボディを評価する．

Egisonの\lstinline{matchAll}式の特徴は，（1）\lstinline{matchAll}という名前と結果としてリストを返すことと，（2）マッチャーという追加の引数をとることにある．
（1）の特徴は，複数の結果をもつパターンマッチをサポートするための特徴である．
\lstinline{matchAll}式は複数のパターンマッチの結果すべてについてボディを評価して，その結果を集めてリストとして返す．
上記の例のパターンに使われている\lstinline{::}\index{::@\lstinline{::}}はコンスパターンと呼ばれるリストを先頭の要素と残りのリストに分解するパターンコンストラクタである．
そのため，上記の例の場合は，パターンマッチの結果が一つであるので，長さが$1$のリストを返している．
（2） の特徴は，パターンマッチアルゴリズムの拡張性とパターンの多相性を実現するための特徴である．
マッチャー\index{まっちゃー@マッチャー}はEgison以外の言語ではみられない独自のオブジェクトで，パターンマッチアルゴリズムを保持するためのオブジェクトである．
マッチャーによるパターンの多相性については，\ref{quick-matcher}節で扱う．

\lstinline{--}\index{--@\lstinline{--}}はEgisonのインライン・コメント・デリミタである．
本書を通して，プログラム直後のインライン・コメントを使って，プログラムの実行結果を記述する．

\begin{lstlisting}[language=egison]
matchAll [1,2,3] as list something with
| $hs ++ $ts -> (hs, ts)
-- [([], [1, 2, 3]), ([1], [2, 3]), ([1, 2], [3]), ([1, 2, 3], [])]
\end{lstlisting}

\begin{lstlisting}[language=egison]
matchAll ターゲット as マッチャー with
| パターン -> ボディ
| パターン -> ボディ
...
\end{lstlisting}

マッチ節が1つであるときは，下記のようにマッチ節の前の\lstinline{|}を省略することができる．

\begin{lstlisting}[language=egison]
matchAll ターゲット as マッチャー with パターン -> ボディ
\end{lstlisting}

\section{値パターンと述語パターンによる非線形パターンの表現}

\lstinline{matchAll}式は，非線形パターンと組み合わせたときにその本領を発揮する．
たとえば，以下の非線形パターンは，対象のコレクションが同値な要素のペアを含む場合にパターンマッチに成功する．

\begin{lstlisting}[language=egison]
matchAll [1,2,3,2,4,3] as list integer with
| _ ++ $x :: _ ++ #x :: _ -> x
-- [2,3]
\end{lstlisting}

\emph{値パターン}\index{あたいぱたーん@値パターン}は非線形パターンを表現するために重要な役割を果たす．
値パターンは，値パターンの中身とターゲットが等しいかどうかチェックする．
値パターンの先頭には，\lstinline{#}\index{#@\lstinline{#}}が付加される．
\lstinline{#}の後ろには任意の式を書くことができる．
\lstinline{#}に続く式は，その値パターンの左側に現れたパターン変数に束縛された値を参照しながら評価される．
この動作を実現するため，Egisonのパターンは左から右に順番に処理されることが決まっている．
その結果，\lstinline{$x :: #x :: _}のようなパターンは妥当であるが，\lstinline{#x :: $x :: _}は間違いである．

よりおもしろい非線形パターンの例として，双子素数のパターンマッチを紹介する．
双子素数とは，$(p, p+2)$という形の素数のペアのことをいう．
\lstinline{primes}には素数の無限列が束縛されている．
\lstinline{primes}はEgisonの組み込みライブラリで定義されている．
この\lstinline{matchAll}式は，素数の無限列からすべての双子素数を順番に抜き出す．

\begin{lstlisting}[language=egison]
twinPrimes := matchAll primes as list integer with
  | _ ++ $p :: #(p + 2) :: _ -> (p, p + 2)

take 8 twinPrimes
-- [(3, 5), (5, 7), (11, 13), (17, 19), (29, 31), (41, 43), (59, 61), (71, 73)]
\end{lstlisting}

パターンマッチのときに，同値性よりも一般的な述語を使いたいことがある．
\emph{述語パターン}\index{じゅつごぱたーん@述語パターン}はそのために用意されている組み込みパターンである．
述語パターンは，述語パターンの中身の述語にターゲットを適用したときに真が帰ってきた場合，パターンマッチに成功するパターンである．
述語パターンの先頭は，\lstinline{?}からはじまり，\lstinline{?}のあとには1引数の述語が続く．

\begin{lstlisting}[language=egison]
twinPrimes = matchAll primes as list integer with
  | _ ++ $p : ?(\q -> q == p + 2) : _ -> (p, p + 2)
\end{lstlisting}

\section{バックトラッキングによる効率的な非線形パターンマッチ}

The pattern-matching algorithm inside Egison includes the backtracking mechanism for efficient non-linear pattern matching.\footnote{The Author's paper~\cite{egi2019macro} on Scheme macros that implement our pattern-matching system discusses the performance of a program in pattern-match-oriented style. A program in pattern-match-oriented style is currently two times slower when compared with the same program in functional programming style.}

\begin{lstlisting}[language=egison]
matchAll [1..n] as list integer with _ ++ $x :: _ ++ #x :: _ -> x
-- O(n^2) で [] を返す
matchAll [1..n] as list integer with _ ++ $x :: _ ++ #x :: _ ++ #x :: _ -> x
-- O(n^2) で [] を返す
\end{lstlisting}

The above expressions match a collection that consists of integers from $1$ to $n$ as a list of integers for enumerating identical pairs and triples, respectively.
This target collection contains neither identical pairs nor triples, therefore both expressions return an empty collection.

When evaluating the second expression, Egison interpreter does not try pattern matching for the second \lstinline|#x| because pattern matching for the first \lstinline|#x| always fails.
Therefore, the time complexities of the above expressions are identical.
The pattern-matching algorithm inside Egison is discussed in~\cite{egi2018Aplas} in detail.

\section{マッチャーによるパターンの拡張性と多相性}\label{quick-matcher}

Another merit of matchers, in addition to the extensibility of pattern-matching algorithms, is the ad-hoc polymorphism of patterns.
The ad-hoc polymorphism of patterns is important for non-free data types because some data are pattern-matched as various non-free data types at the different parts of a program.
For example, a collection is pattern-matched as a list, a multiset, and a set.
Polymorphic patterns reduce the number of names for pattern constructors.

In the following sample, a list \lstinline{[1,2,3]} is pattern-matched using different matchers with the same cons pattern.
In the case of sets, the rest elements are the same as the original collection because we ignore the redundant elements.
If we regard a set as a collection that contains infinitely many copies of each element, this specification of the cons pattern for sets is natural.

\begin{lstlisting}[language=egison]
matchAll [1,2,3] as list something with $x :: $xs -> (x,xs)
-- [(1,[2,3])]
matchAll [1,2,3] as multiset something with $x :: $xs -> (x,xs)
-- [(1,[2,3]),(2,[1,3]),(3,[1,2])]
matchAll [1,2,3] as set something with $x :: $xs -> (x,xs)
-- [(1,[1,2,3]),(2,[1,2,3]),(3,[1,2,3])]
\end{lstlisting}

Polymorphic patterns are useful especially for value patterns.
As well as other patterns, the behavior of value patterns is dependent on matchers.
For example, an equality \lstinline{[1,2,3] == [2,1,3]} between collections is false if we regard them as lists but true if we regard them as multisets.
Still, thanks to ad-hoc polymorphism of patterns, we can use the same syntax for both types.
This dramatically improves the readability of the program and makes programming with non-free data types easy.

\begin{lstlisting}[language=egison]
matchAll [1,2,3] as list integer with #[2,1,3] -> "Matched" -- []
matchAll [1,2,3] as multiset integer with #[2,1,3] -> "Matched" -- ["Matched"]
\end{lstlisting}

\section{\lstinline{matchAllDFS}式によるパターンマッチ結果の順序の制御}

The \lstinline{matchAll} expression is designed to enumerate all countably infinite pattern-matching results.
For this purpose, users sometimes need to care about the order of pattern-matching results.

Let us start by showing a representative sample.
The \lstinline{matchAll} expression below enumerates all pairs of natural numbers.
We extract the first $8$ elements with the \lstinline{take} function.
\lstinline{matchAll} traverses the reduction tree of pattern matching in breadth-first search to traverse all the nodes~(Sect. 5.2 of~\cite{egi2018Aplas}).
As a result, the order of the pattern-matching results is as follows.

\begin{lstlisting}[language=egison]
take 8 (matchAll [1..] as set something with $x : $y : _ -> (x,y))
-- [(1,1),(1,2),(2,1),(1,3),(2,2),(3,1),(2,3),(3,2)]
\end{lstlisting}

The above order is preferable for traversing an infinitely large reduction tree.
However, sometimes, this order is not preferable (see Sect.~\ref{nested-join-cons} and~\ref{trees}).
\lstinline{matchAllDFS}\index{matchAllDFS@\lstinline{matchAllDFS}} that traverses a reduction tree in depth-first order is provided for this reason.

\begin{lstlisting}[language=egison]
take 8 (matchAllDFS [1..] as set something with $x : $y : _ -> (x,y))
-- [(1,1),(1,2),(1,3),(1,4),(1,5),(1,6),(1,7),(1,8)]
\end{lstlisting}

\section{andパターン・orパターン・notパターン}

The situations where and-patterns and or-patterns are useful are similar to those of the existing languages, whereas not-patterns become useful when they are combined with non-linear pattern matching with backtracking.

We start by showing pattern matching for \emph{prime triples} as an example of and-patterns and or-patterns.
A prime triple is a triple of primes whose form is $(p,p+2,p+6)$ or $(p,p+4,p+6)$.
The and-pattern is used as an as-pattern.
The or-pattern is used to match both of $p+2$ and $p+4$.

\begin{lstlisting}[language=egison]
primeTriples = matchAll primes as list integer with
                  _ ++ $p : (and (or #(p + 2) #(p + 4)) $m) : #(p + 6) : _ -> (p, m, p + 6)

take 8 primeTriples -- [(5,7,11),(7,11,13),(11,13,17),(13,17,19),(17,19,23),(37,41,43),(41,43,47),(67,71,73)]
\end{lstlisting}


A not-pattern matches a target if the pattern does not match the target, as its name implies.
A not-pattern is prepended with \lstinline{!}, and a pattern follows after \lstinline{!}.
The following \lstinline{matchAll} enumerates sequential pairs of prime numbers that are \emph{not} twin primes.

\begin{lstlisting}[language=egison]
take 10 (matchAll primes as list integer with _ ++ $p : (and !#(p + 2) $q) : _ -> (p, q))
-- [(2,3),(7,11),(13,17),(19,23),(23,29),(31,37),(37,41),(43,47),(47,53),(53,59)]
\end{lstlisting}

\section{ループパターン}

A loop pattern is a pattern construct for representing a pattern that repeats multiple times.
It is an extension of Kleene star operator of regular expressions for general non-free data types~\cite{egi2018loop}.

Let us start by considering pattern matching for enumerating all combinations of two elements from a target collection.
It can be written using \lstinline{matchAll} as follows.

\begin{lstlisting}[language=egison]
comb2 xs = matchAll xs as list something with _ ++ $x_1 : _ ++ $x_2 : _ -> [x_1, x_2]

comb2 [1,2,3,4] -- [[1,2],[1,3],[2,3],[1,4],[2,4],[3,4]]
\end{lstlisting}

Egison allows users to append indices to a pattern variable as \lstinline|$x_1| and \lstinline|$x_2| in the above sample.
They are called \emph{indexed variables} and represent $x_1$ and $x_2$ in mathematical expressions.
The expression after \lstinline|_| must be evaluated to an integer and is called an \textit{index}.
We can append as many indices as we want like \lstinline|x_i_j_k|.
When a value is bound to an indexed pattern variable \lstinline|$x_i|, the system initiates an abstract map consisting of key-value pairs if \lstinline|x| is not bound to a map, and bind it to \lstinline{x}.
If \lstinline|x| is already bound to a map, a new key-value pair is added to this map.

Now, we generalize \lstinline{comb2}.
The loop patterns can be used for that purpose.

\begin{lstlisting}[language=egison]
comb n xs = matchAll xs as list something with
              loop $i (1,n)
                (_ ++ $x_i : ...)
                _ -> map (\i -> x_i) [1..n]

comb 2 [1,2,3,4] -- [[1,2],[1,3],[2,3],[1,4],[2,4],[3,4]]
comb 3 [1,2,3,4] -- [[1,2,3],[1,2,4],[1,3,4],[2,3,4]]
\end{lstlisting}

The loop pattern takes an \emph{index variable}, \emph{index range}, \emph{repeat pattern}, and \emph{end pattern} as arguments.
An index variable is a variable to hold the current repeat count.
An index range specifies the range where the index variable moves.
A repeat pattern is a pattern repeated when the index variable is in the index range.
An end pattern is a pattern expanded when the index variable gets out of the index range.

Inside loop patterns, we can use the \emph{ellipsis pattern} (\lstinline{...}).
The repeat pattern or the end pattern is expanded at the location of the ellipsis pattern.
The repeat pattern is expanded replacing the ellipsis pattern incrementing the value of the index variable.

%The important difference between loop patterns and recursive patterns~\cite{queinnec1990compilation} is that loop patterns provide a simple method for changing the content of the pattern repeated depending on the repeat count.
%Users can change the content of the repeat pattern by referring to the repeat count through the index variable.
%Furthermore, indexed variables that often refer to the index variable in their index (such as \lstinline|$x_i| in the above samples) enable us to describe more expressive patterns by allowing us to refer to the values bound to the indexed pattern variable in the pattern previously repeated.

\medskip

The repeat counts of the loop patterns in the above samples are constants.
However, we can also write a loop pattern whose repeat count varies depending on the target by specifying a pattern instead of an integer as the end number.
When an end number is a pattern, the ellipsis pattern is replaced with both the repeat pattern and the end pattern, and the repeat count when the ellipsis pattern is replaced with the end pattern is pattern-matched with that pattern.
The following loop pattern enumerates all initial prefixes of the target collection.

{\footnotesize
\begin{lstlisting}[language=egison]
matchAll [1,2,3,4] as list something with loop $i (1, $n) ($x_i : ...) _ -> map (\i -> x_i) [1..n]
-- [[],[1],[1,2],[1,2,3],[1,2,3,4]]
\end{lstlisting}
}

Loop patterns are heavily used especially for trees and graphs.
We work on pattern matching for trees in Sect.~\ref{trees}.
%Pattern matching for graphs is demonstrated in Appendix~\ref{graphs}.
More formal specification of syntax and semantics of loop patterns is shown in the author's previous paper~\cite{egi2018loop}.

\section{シーケンシャルパターン}

The pattern-matching system of Egison processes patterns from left to right in order.
However, there are cases where we want to change this order, for example, to refer to the value bound to the right side of a pattern.
Sequential patterns are provided for such a purpose.


Sequential patterns allow users to control the order of the pattern-matching process.
A sequential pattern is represented as a list of patterns.
Pattern matching is executed for each pattern in order.
In the following sample, the target list is pattern-matched from the third, first, and second element in order.

\begin{lstlisting}[language=egison]
matchAll [2,3,1,4,5] as list integer with
  [ @ :       @       : $x : _,
   (#(x + 1), @),
   #(x + 2)] -> "Matched" -- ["Matched"]
\end{lstlisting}

\noindent
\lstinline|@| that appears in a sequential pattern is called \emph{later pattern variable}.
The target data bound to later pattern variables are pattern-matched in the next sequence.
When multiple later pattern variables appear, they are pattern-matched as a tuple in the next sequence.
It allows us to apply not-patterns for different parts of a pattern at the same time as we will see in Sect~\ref{tuple-of-collections}.

Some readers might wonder that a sequential pattern can be transformed into a nested \lstinline{match-all} expression.
There are at least two reasons why it is impossible.
First, a nested \lstinline{match-all} expression breaks breadth-first search strategy: the inner \lstinline{match-all} for the second result of the outer \lstinline{match-all} is executed only after the inner \lstinline{match-all} for the first result of the outer \lstinline{match-all} is finished.
Second, a later pattern variable retains the information of not only a target but also a matcher.
There are cases that the matcher of \lstinline{match-all} is a parameter passed as an argument of a function, and a pattern is polymorphic.
Therefore, it is impossible to determine the matchers of inner \lstinline{match-all} expressions syntactically.

\section{forallパターン}

述語パターンで奇数かどうかチェックしている．
パターンマッチに成功した場合，すべてのパターンマッチ結果を返す．

\begin{lstlisting}[language=egison]
matchAll [1,5,3] as multiset integer with
| forall ((@ & $x) :: _) ?odd? -> x
-- [1,5,3]
\end{lstlisting}

1つでも第二引数のパターンのパターンマッチに失敗する第一引数のパターンマッチの結果がある場合，全体のパターンマッチに失敗する．

\begin{lstlisting}[language=egison]
matchAll [1,4,3] as multiset integer with
| forall ((@ & $x) :: _) ?odd? -> x
-- []
\end{lstlisting}

notパターンでもforallパターンに近い意味のパターンを表現できる．
notパターンのforallパターンに対する利点は，以下の二点である．
\begin{itemize}
\item[(1)] forallパターン中のパターン変数に値を束縛し，複数のパターンマッチの結果すべてを返すことができる（notパターンの中身ではパターン変数に値を束縛することができない）．
\item[(2)] forallパターンを使ったほうがより直接的に意味を表現できる場合がある．
\end{itemize}

\begin{lstlisting}[language=egison]
matchAll [1,5,3] as multiset integer with
| !(!?odd? :: _) -> ``match''
-- [``match'']
\end{lstlisting}

\section{パターン関数によるパターンのモジュール化}


\section{マッチャー合成による新しいマッチャーの生成}

Matchers are composable, and we can define matchers for such as tuples of multisets and multisets of multisets.
Using this feature, we can define matchers for various data types.

First, we can define a matcher for tuples by a tuple of matchers.
A tuple pattern is used for pattern matching using such a matcher.
For example, we can define the \lstinline{intersect} function using a matcher for tuples of two multisets.
We work on pattern matching for tuples of collections more in Sect.~\ref{tuple-of-collections}.

\begin{lstlisting}[language=egison]
intersect xs ys = matchAll (xs,ys) as (multiset eq, multiset eq) with ($x : _, #x : _) -> x
\end{lstlisting}

\noindent \lstinline{eq} is a user-defined matcher for data types for which equality is defined. 
When the \lstinline{eq} matcher is used, equality is checked for a value pattern.\footnote{A definition of the \lstinline{eq} matcher is explained in Sect.~6.3 of~\cite{egi2018Aplas}.}
%In fact, the \lstinline{integer} matcher used in the previous sections is an alias of the \lstinline{eq} matcher.

By passing a tuple matcher to a function that takes and returns a matcher, we can define a matcher for various non-free data types.
For example, we can define a matcher for a graph as a set of edges.
In the following code, we assume a node id is represented by an integer.

{\footnotesize
\begin{lstlisting}[language=egison]
graph = multiset (integer, integer)
\end{lstlisting}
}

\noindent A matcher for adjacency graphs also can be defined.
An adjacency graph is defined as a multiset of tuples of an integer and a multiset of integers.

{\footnotesize
\begin{lstlisting}[language=egison]
adjacencyGraph = multiset (integer, multiset integer)
\end{lstlisting}
}

%\noindent We demonstrate pattern matching for these graphs in Appendix~\ref{graphs}.

\medskip

Some readers might wonder about matchers for algebraic data types.
Egison provides a special syntax construct for defining a matcher for an algebraic data type.
For example, a matcher for binary trees can be defined using \lstinline{algebraicDataMatcher}.

{\footnotesize
\begin{lstlisting}[language=egison]
algebraicDataMatcher binaryTree a = BLeaf a | BNode a (binaryTree a) (binaryTree a)
\end{lstlisting}
}

Matchers for algebraic data types and matchers for non-free data types also can be combined.
For example, we can define a matcher for trees whose nodes have an arbitrary number of children whose order is ignorable.
We show pattern matching for these trees in Sect.~\ref{trees}.

{\footnotesize
\begin{lstlisting}[language=egison]
algebraicDataMatcher tree a = Leaf a | Node a (multiset (tree a))
\end{lstlisting}
}
